\chapter{Thấy mình để sửa}
\markboth{Thấy mình để sửa}{See Yourself and Change}

\section{Một người chỉ thay đổi khi nghe chính câu nói của mình từ miệng người khác.}
\begin{truyen}{Một người chỉ thay đổi khi nghe chính câu nói của mình từ miệng người khác.}{Sometimes mirrors are other people.}
\chuhoa{K}{hi nghe người khác lặp lại đúng giọng điệu mình từng dùng, anh mới giật mình vì sự lạnh lùng trong đó. Điều ta quen ở mình nhiều khi phải đi một vòng qua người khác mới nhìn ra được.}
\textit{(Hearing someone else repeat his own tone shocked him with its harshness. What we are used to in ourselves often needs to return through others before we can see it.)}
\end{truyen}
\begin{baihoc}
Nhìn ra mình là khởi đầu của sửa mình; không có bước đó thì mọi lý thuyết đúng sai đều rất mỏng.
\textit{(Seeing yourself is the beginning of changing yourself; without that step, all theories of right and wrong stay thin.)}
\end{baihoc}
\begin{ghinhoanh}\textbf{Short English to remember:}

Sometimes mirrors are other people.
\end{ghinhoanh}
\begin{tuhocnhanh}
1. Có phản ứng nào của người khác đang soi lại chính em không? \textit{(Is someone else's reaction mirroring you back right now?)}

2. Em nhìn ra điều gì về mình từ đó? \textit{(What do you see about yourself from it?)}
\end{tuhocnhanh}
\ngancach

\section{Một người từng hay đổ lỗi rồi quyết dừng lại.}
\begin{truyen}{Một người từng hay đổ lỗi rồi quyết dừng lại.}{Owning your part changes the story.}
\chuhoa{T}{rước đây chuyện gì anh cũng tìm được ai đó để trách. Từ khi tập nhận phần của mình, anh bớt cay cú và tiến bộ nhanh hơn. Không phải vì đời bỗng dễ hơn, mà vì anh ngừng lãng phí năng lượng vào việc trốn tránh.}
\textit{(He used to find someone to blame in every situation. Once he started owning his part, he became less bitter and improved faster. Life did not get easier; he simply stopped wasting energy avoiding responsibility.)}
\end{truyen}
\begin{baihoc}
Nhận phần mình không làm gánh nặng hơn; nó làm ta có quyền sửa hơn.
\textit{(Owning your part does not make the burden heavier; it gives you more power to change it.)}
\end{baihoc}
\begin{ghinhoanh}\textbf{Short English to remember:}

Owning your part changes the story.
\end{ghinhoanh}
\begin{tuhocnhanh}
1. Trong chuyện gần nhất, phần của em là gì? \textit{(In the most recent conflict, what was your part?)}

2. Em có đang né trách nhiệm bằng cách nào? \textit{(How are you avoiding responsibility?)}
\end{tuhocnhanh}
\ngancach

\section{Một người biết mình sai nhưng sửa rất chậm vì sĩ diện.}
\begin{truyen}{Một người biết mình sai nhưng sửa rất chậm vì sĩ diện.}{Pride slows repentance.}
\chuhoa{A}{nh thấy rõ mình cần đổi, nhưng cái tôi khiến anh kéo dài thêm nhiều tháng. Mỗi ngày chậm sửa là mỗi ngày để cái sai ăn sâu hơn. Thấy mình mà không chịu sửa nhanh cũng là một kiểu tự làm khó đời mình.}
\textit{(He clearly saw he needed to change, but pride made him delay for months. Every day of delay let the wrong habit sink deeper. Seeing yourself without changing quickly is another way of making life harder.)}
\end{truyen}
\begin{baihoc}
Điều khó không phải lúc nào cũng là nhìn ra lỗi; nhiều khi là hạ cái tôi xuống để sửa ngay.
\textit{(The hardest part is not always seeing the fault; sometimes it is lowering ego enough to fix it immediately.)}
\end{baihoc}
\begin{ghinhoanh}\textbf{Short English to remember:}

Pride slows repentance.
\end{ghinhoanh}
\begin{tuhocnhanh}
1. Lỗi nào em đã nhìn ra nhưng còn trì hoãn sửa? \textit{(What fault have you recognized but still delayed fixing?)}

2. Nếu sửa ngay hôm nay, bước đầu là gì? \textit{(If you fixed it today, what would the first step be?)}
\end{tuhocnhanh}
\ngancach

\section{Một cô gái thay đổi chỉ sau khi thôi nói mình là như vậy rồi.}
\begin{truyen}{Một cô gái thay đổi chỉ sau khi thôi nói mình là như vậy rồi.}{Identity should not excuse habit.}
\chuhoa{C}{ô từng dùng tính cách làm lá chắn: em nóng tính vậy đó, em nói thẳng vậy đó. Khi thôi lấy bản thân hiện tại làm lý do bất biến, cô mới bắt đầu lớn lên thật.}
\textit{(She used personality as a shield: I am just hot-tempered, I just speak bluntly. Once she stopped treating her current self as a permanent excuse, real growth began.)}
\end{truyen}
\begin{baihoc}
Tính cách giải thích khuynh hướng, nhưng không nên trở thành nơi trú ẩn cho thói xấu.
\textit{(Personality may explain tendencies, but it should not become shelter for bad habits.)}
\end{baihoc}
\begin{ghinhoanh}\textbf{Short English to remember:}

Identity should not excuse habit.
\end{ghinhoanh}
\begin{tuhocnhanh}
1. Em có đang dùng câu đó là con người em để khỏi sửa không? \textit{(Are you using that's just who I am to avoid change?)}

2. Thói quen nào cần được tách khỏi bản sắc của em? \textit{(What habit needs to be separated from your identity?)}
\end{tuhocnhanh}
\ngancach

\section{Một người về già chỉ quý nhất khả năng tự sửa mình sớm.}
\begin{truyen}{Một người về già chỉ quý nhất khả năng tự sửa mình sớm.}{The quickest repair is a hidden wisdom.}
\chuhoa{N}{hìn lại đời mình, ông không thấy người giỏi nhất là người chưa từng sai, mà là người sai vừa đủ để hiểu rồi sửa sớm. Chính tốc độ quay đầu ấy cứu họ khỏi nhiều năm sống lệch.}
\textit{(Looking back on life, he did not see the wisest people as those who never failed, but those who failed enough to understand and corrected themselves early. That speed of turning back saved them from many years of living crookedly.)}
\end{truyen}
\begin{baihoc}
Khôn ngoan không phải không bao giờ lạc; khôn ngoan là lạc rồi quay lại sớm.
\textit{(Wisdom is not never getting lost; wisdom is returning early after getting lost.)}
\end{baihoc}
\begin{ghinhoanh}\textbf{Short English to remember:}

The quickest repair is hidden wisdom.
\end{ghinhoanh}
\begin{tuhocnhanh}
1. Em muốn quay lại điều gì sớm hơn trước khi quá muộn? \textit{(What do you want to return to sooner before it becomes too late?)}

2. Từ nay, em muốn sửa mình nhanh hơn ở điểm nào? \textit{(From now on, where do you want to correct yourself faster?)}
\end{tuhocnhanh}
\ngancach